\documentclass[11pt]{article}
\usepackage[margin=1in]{geometry}
\usepackage{amsmath, amssymb, amsthm}
\usepackage{mathtools}
\usepackage{hyperref}

\newtheorem{theorem}{Theorem}
\newtheorem{lemma}[theorem]{Lemma}
\newtheorem{proposition}[theorem]{Proposition}
\newtheorem{corollary}[theorem]{Corollary}
\newtheorem{conjecture}[theorem]{Conjecture}
\theoremstyle{definition}
\newtheorem{definition}[theorem]{Definition}
\newtheorem{remark}[theorem]{Remark}

\newcommand{\Sym}{\mathfrak{S}}
\newcommand{\Hq}{\mathcal{H}_q}
\newcommand{\la}{\lambda}
\newcommand{\hla}{\hat{\lambda}}
\newcommand{\Om}{\Omega}
\newcommand{\Omstar}{\Omega^{*}}
\newcommand{\vT}{v_T}
\newcommand{\ZN}{\mathbb{Z}_{\ge 0}}

\title{The SC-at-$1$ trivial reduction\\
of Converse Diagnostic 4}
\author{Clio}
\date{2026-05-19}

\begin{document}
\maketitle

\begin{abstract}
We give a trivial proof closing $\mathbf{378/514 \approx 73.5\%}$ of
the empirically tested converse Diagnostic 4 cases. For
$T \in \mathrm{SYT}(\la)$ with $\la = (\hla, 1^q)$ and $\tau(\hla) = 0$,
if the cell containing $2$ in $T$ is directly below the cell containing
$1$ (an SC pair at index $1$), then $S_1 \vT = 0$ in the Hoefsmit
basis of $V^\la$. Hence $R'_n \vT = 0$, and thus
$\Om \vT = 0$ and $p_T(q) = 0$. Combined with yesterday's universal SC
necessity proof, this fully closes Diagnostic 4 (forward and converse)
for the SC-at-$1$ stratum. The conjecture's remaining open core is the
$\mathrm{SR}$-at-$1$ stratum: $136$ zero diagonals out of $514$, all of
which $\Om$- or $\Omstar$-kernel-vanish empirically, but for which the
chain-product collapses at deeper combinatorial events.
\end{abstract}

\section{Setting}

Fix $\la = (\hla, 1^q)$ with $\ell(\hla) \ge 2$, every row
$\hla_r \ge 2$, and $\tau(\hla) := \max(0, q + 2\ell(\hla) - 1 - |\hla|)
= 0$. Let $n = |\la|$, $\ell = \ell(\la)$, and let $V^\la$ denote the
irreducible $\Hq(\Sym_n)$-module realised in the Hoefsmit (seminormal)
basis indexed by $\mathrm{SYT}(\la)$. We write $T_i$ for the Hecke
generators with $(T_i - q)(T_i + 1) = 0$ and put $S_i := T_i + 1$.

Following the conventions of yesterday's note
\texttt{2026-05-18-converse-diagnostic-4.tex}, define
\begin{align*}
R'_k &:= S_{k-1} S_{k-2} \cdots S_2 S_1, \\
L_k  &:= S_1 S_2 \cdots S_{k-2} S_{k-1}, \\
\Om   &:= R'_{\ell+1} R'_{\ell+2} \cdots R'_n, \\
\Omstar &:= L_n L_{n-1} \cdots L_{\ell+1},
\end{align*}
and write $\vT \in V^\la$ for the seminormal basis vector indexed by
$T$. The diagonal probe is
\[
p_T(q) := \bra{\vT}\Om\ket{\vT} \in \mathbb{Q}(q),
\]
where $\bra{\cdot}\cdot\ket{\cdot}$ pairs a basis vector with the
coefficient of $\vT$ in the expansion.

\begin{definition}
An \emph{SC pair} of $T$ at index $i$ is a pair $(i, i+1)$ such that
the cells containing $i$ and $i+1$ in $T$ share a column. An
\emph{SR pair} is the row-analogue. We say $T$ is \emph{SC-at-$1$} if
the cell containing $2$ is $(2, 1)$ (so $2$ is directly below $1$);
equivalently, the SC pair $(1, 2)$ is present. Otherwise (since $2$ is
either at $(1, 2)$ or $(2, 1)$ in any SYT) we say $T$ is
\emph{SR-at-$1$}.
\end{definition}

\section{The trivial step}

\begin{lemma}\label{lem:S1-kills}
For $T \in \mathrm{SYT}(\la)$, if $T$ is SC-at-$1$ then
$S_1 \vT = 0$.
\end{lemma}

\begin{proof}
By Hoefsmit's seminormal formula, if entries $i, i+1$ of $T$ share a
column, then $T_i \vT = -\vT$. Applying this at $i = 1$ for an SC-at-$1$
$T$ gives $T_1 \vT = -\vT$, whence $S_1 \vT = (T_1 + 1)\vT = 0$.
\end{proof}

\begin{lemma}\label{lem:Rn-kills}
For any vector $w \in V^\la$ with $S_1 w = 0$ we have $R'_n w = 0$. In
particular, if $T$ is SC-at-$1$ then $R'_n \vT = 0$.
\end{lemma}

\begin{proof}
By definition $R'_n = S_{n-1} S_{n-2} \cdots S_2 S_1$, whose rightmost
factor is $S_1$. Hence
$R'_n w = S_{n-1} \cdots S_2 (S_1 w) = S_{n-1} \cdots S_2 \cdot 0 = 0$.
Combined with Lemma~\ref{lem:S1-kills}, $S_1 \vT = 0$, so
$R'_n \vT = 0$.
\end{proof}

\begin{theorem}[SC-at-$1$ trivial reduction]\label{thm:main}
Let $\la$ be a $\tau = 0$ shape as above. For every
$T \in \mathrm{SYT}(\la)$ with $T$ SC-at-$1$,
\[
\Om \vT = 0
\quad\text{and consequently}\quad
p_T(q) = 0.
\]
\end{theorem}

\begin{proof}
By Lemma~\ref{lem:Rn-kills}, $R'_n \vT = 0$. Then
\[
\Om \vT = R'_{\ell+1} R'_{\ell+2} \cdots R'_{n-1} (R'_n \vT)
= R'_{\ell+1} \cdots R'_{n-1} \cdot 0 = 0,
\]
so $\Om \vT = 0$ and the diagonal entry
$p_T(q) = \bra{\vT}\Om\ket{\vT} = 0$.
\end{proof}

\begin{remark}
Theorem~\ref{thm:main} is striking only against the backdrop of the
forward direction. Yesterday's note showed
$p_T = 0 \Rightarrow T \text{ has some SC pair}$, but the SC pair could
be at any index. Theorem~\ref{thm:main} says that if it is at the
\emph{specific} index $1$, then converse Diagnostic 4 closes by a
one-line algebraic identity: the Hecke factor $S_1$ at the right
end of $R'_n$ already kills $\vT$. No further argument is needed.
\end{remark}

\section{Computational check}

In a sample of $12$ shapes
\[
\la \in \{(3,2),\,(3,2,1),\,(2,2,1),\,(3,2,2),\,(3,2,1,1),\,(3,3,2),\,
(3,2,2,1),\,(3,3,2,1),\,(4,3),\,(4,3,1),\,(4,3,2),\,(3,3,1,1,1)\}
\]
covering every previously-tested converse target (cf.\
\texttt{2026-05-18-converse-probe/dichotomy\_verifier.py}), we
symbolically verified $S_1 \vT = 0$ for every $T$ with $2$ in cell
$(2,1)$ in the Hoefsmit basis. The count was $378$ out of $378$ such
$T$ (i.e.\ no failures). Cross-checked against the stratification
trace
\texttt{2026-05-19-stratification/stratify.py}, all $378$ SC-at-$1$
$T$ produced the collapse step $(k_R, i_R) = (n, 1)$, meaning the very
first factor $S_1$ of $R'_n$ vanishes the basis vector.

This is a confirmation of Lemma~\ref{lem:S1-kills} on actual matrices,
not a separate piece of evidence: Hoefsmit's formula is exact and the
identity $S_1 \vT = 0$ is a direct calculation. We report it because
matching the predicted collapse step at the symbolic level is the
cleanest external check.

\section{Consequence: the residual open problem}

Combining Theorem~\ref{thm:main} with yesterday's universal SC
necessity (every $T$ with $p_T = 0$ has at least one SC pair) gives:

\begin{corollary}\label{cor:residual}
For $\la$ of the form above, the set of zero diagonals decomposes as
\[
\{T \in \mathrm{SYT}(\la) : p_T = 0\}
= \{T : T \text{ is SC-at-}1\}
\;\sqcup\;
\{T : T \text{ is SR-at-}1 \text{ and } p_T = 0\}.
\]
The first set is closed under converse Diagnostic 4 by
Theorem~\ref{thm:main}: every such $T$ satisfies $\Om \vT = 0$.

The full converse conjecture therefore reduces to:
\begin{quote}
For every $T \in \mathrm{SYT}(\la)$ with $T$ SR-at-$1$ (i.e.\ $2 \in
(1,2)$) and $p_T(q) = 0$, either $\Om \vT = 0$ or $\Omstar \vT = 0$.
\end{quote}
\end{corollary}

In our $12$-shape corpus the SC-at-$1$ stratum contributes
$378/514 \approx 73.5\%$ of all zero diagonals. The SR-at-$1$ residual
contains $136$ zero diagonals.

\section{Empirical scaffolding for the residual}

Although the SR-at-$1$ converse is not closed in this note, the
stratification trace produces tight constraints. We summarise what
the data force.

\paragraph{Right-collapse normal form.}
For every SR-at-$1$ zero diagonal in the corpus, the collapse step
$(k_R, i_R)$ of the right chain $\Om$, if finite, has the form
$(k_R, 1)$: \emph{the leftmost factor $S_1$ of $R'_{k_R}$ vanishes the
running vector}. Equivalently, the partial product
$R'_{k_R + 1} R'_{k_R + 2} \cdots R'_n \vT$ already lies in
$\ker S_1$, i.e.\ the eigenspace where $T_1$ acts by $-1$:
\[
\mathrm{span}\{v_S : S \in \mathrm{SYT}(\la),\; S \text{ is SC-at-}1\}.
\]

\paragraph{Left-collapse normal form.}
For every SR-at-$1$ zero diagonal whose $\Omstar$-collapse is finite,
$(k_L, i_L) = (k_L,\, k_L - 1)$: \emph{the rightmost factor
$S_{k_L - 1}$ of $L_{k_L}$ vanishes the running vector}. Equivalently
the partial product $L_{k_L - 1} \cdots L_{\ell + 1} \vT$ lies in
$\ker S_{k_L - 1}$.

\paragraph{$\Om$-survivors.}
$12$ of the $136$ SR-at-$1$ zeros have $\Om \vT \ne 0$ (and are killed
only by $\Omstar$). Every one of them has either
\[
\text{(i) } T \text{ contains } 1, 2, 3 \text{ in row } 1
\qquad\text{(11 of 12),}
\]
\[
\text{(ii) } T \text{ has } 1, 2 \text{ in row } 1 \text{ and } 3, 4
\text{ in row } 2 \text{ at cells }(2,1),(2,2) \text{ (1 of 12).}
\]
In both cases $T$ has at least \emph{two} SR pairs at small indices,
preventing the chain from steering into $\ker S_1$ during~$\Om$.

This phenomenology suggests the residual conjecture splits into a
\emph{double-SR-prefix} sub-case (handled by $\Omstar$) and a
\emph{single-SR-prefix} sub-case (handled by $\Om$).

\begin{conjecture}[Stratified converse Diagnostic 4]\label{conj:strat}
Let $T \in \mathrm{SYT}(\la)$ with $T$ SR-at-$1$ and $p_T(q) = 0$.
\begin{enumerate}
\item[\textup{(i)}] If $T$ does not contain $1, 2, 3$ in row $1$ and
$3, 4$ do not both lie in row $2$ at $(2,1), (2,2)$, then
$\Om \vT = 0$.
\item[\textup{(ii)}] Otherwise $\Omstar \vT = 0$.
\end{enumerate}
\end{conjecture}

Conjecture~\ref{conj:strat} holds for all $136$ SR-at-$1$ zero
diagonals in the corpus. Its proof remains open.

\section{Where this leaves Route A}

The Hecke-algebraic stratification (Route A) of the prove session's
plan is now \emph{partially closed}:
\begin{itemize}
\item The SC-at-$1$ stratum closes by Theorem~\ref{thm:main}.
\item The SR-at-$1$ stratum reduces to Conjecture~\ref{conj:strat},
      whose two halves each ask for a different structural identity
      in the chain product.
\end{itemize}
The two halves are not symmetric: SR-at-$1$ T's break the left/right
duality that the SC-at-$1$ trivial proof exploits. This matches the
$2026$-$05$-$19$ refutation of universal SR necessity
(\texttt{2026-05-19-sr-conjecture-refuted.tex}).

\paragraph{Sharper next step.}
In the right-collapse normal form, the chain product
$R'_{k_R + 1} \cdots R'_n \vT$ projects $\vT$ into the SC-at-$1$
eigenspace. Identifying the smallest $k$ at which this projection
becomes total --- and matching it to combinatorial data on $T$ --- is
the natural next problem. This is a question about how the
``letter-shift'' operators $R'_k$ act on the position of $2$ in the
support of the running vector.

\section*{Acknowledgements}

This note builds on yesterday's universal SC necessity
(\texttt{2026-05-18-converse-diagnostic-4.tex}), the two-sided kernel
dichotomy (\texttt{2026-05-18-d-class-row-zero.tex}), and the
SR-necessity refutation
(\texttt{2026-05-19-sr-conjecture-refuted.tex}). The Hoefsmit basis
machinery is from
\texttt{2026-05-08-rank-Pi/verify\_seminormal.py}. Stratification
data from
\texttt{2026-05-19-stratification/stratify.py}.

\end{document}
